\chapter{Konstrukcje przestrzeni topologicznych}
\label{ch:topog-konstrukcje}

\noindent\textbf{Cel rozdziału.}
Dotychczas rozpatrywaliśmy gotowe przestrzenie topologiczne i odwzorowania
między nimi. Teraz pokażemy, jak z danych przestrzeni otrzymywać nowe.
Konstrukcje będą określane przez wymaganie ciągłości naturalnych odwzorowań.
Najpierw wprowadzimy topologię podprzestrzeni, następnie iloczynu i sumy
rozłącznej, a na końcu topologie początkowe i końcowe, w szczególności
topologię ilorazową. Wyprowadzimy jawne kryteria otwartości i własności
uniwersalne, nie zakładając metryzowalności ani aksjomatów oddzielania.

\section{Topologia podprzestrzeni}
\label{sec:topog-podprzestrzen}

\begin{definicja}[Topologia podprzestrzeni]
\label{def:topog-podprzestrzen}
Niech $(X,\tau)$ będzie przestrzenią topologiczną i niech $A\subseteq X$.
\emph{Topologią podprzestrzeni} na $A$ nazywamy rodzinę
\[
\tau|_A=\{A\cap U:U\in\tau\}.
\]
Elementy $\tau|_A$ nazywamy zbiorami otwartymi \emph{w $A$}.
Dotyczy to również $A=\varnothing$.
\end{definicja}

\begin{stwierdzenie}[Aksjomaty topologii podprzestrzeni]
\label{prop:topog-podprzestrzen-topologia}
Dla każdej przestrzeni topologicznej $(X,\tau)$ i każdego $A\subseteq X$
rodzina $\tau|_A$ jest topologią na $A$.
\end{stwierdzenie}
\begin{proof}
Z równości $A\cap\varnothing=\varnothing$ i $A\cap X=A$
wynika obecność $\varnothing$ i $A$. Dla dowolnej rodziny $(U_i)_{i\in I}$
zbiorów otwartych w $X$ mamy
\[
\bigcup_{i\in I}(A\cap U_i)=A\cap\bigcup_{i\in I}U_i,
\]
a dla $U,V\in\tau$ zachodzi
$(A\cap U)\cap(A\cap V)=A\cap(U\cap V)$.
Prawa strona każdej równości należy do $\tau|_A$, ponieważ $\tau$
jest topologią. Argument działa także dla pustej rodziny.
\end{proof}

\begin{przyklad}[Otwarty względnie, lecz nie w całej przestrzeni]
\label{ex:topog-otwarty-wzglednie}
Niech $X=\mathbb R$ z topologią euklidesową i $A=[0,2]$.
Zbiór $[0,1)$ jest otwarty w $A$, ponieważ
$[0,1)=A\cap(-1,1)$, ale nie jest otwarty w $\mathbb R$.
Podobnie $\{0\}$ jest otwarty w podprzestrzeni
$\{0\}\cup[2,3]$, chociaż nie jest otwarty w $\mathbb R$.
\end{przyklad}

\begin{figure}[htbp]
\centering
\begin{tikzpicture}[>=Stealth,scale=1]
 \draw[->] (-2.25,0)--(4.15,0) node[right] {$\mathbb R$};
 \foreach \x in {-2,-1,0,1,2,3} \draw (\x,.08)--(\x,-.08);
 \draw[line width=2.8pt,blue!55] (0,0)--(2,0);
 \fill[blue!70!black] (0,0) circle (2.3pt);
 \fill[blue!70!black] (2,0) circle (2.3pt);
 \draw[line width=2.3pt,red!75!black] (0,.34)--(1,.34);
 \fill[red!75!black] (0,.34) circle (2.7pt);
 \draw[fill=white,draw=red!75!black,thick] (1,.34) circle (2.7pt);
 \node[blue!65!black,above] at (2.35,.05) {$A=[0,2]$};
 \node[red!75!black,above] at (.48,.42) {$A\cap(-1,1)=[0,1)$};
 \node[below] at (0,-.11) {$0$};
 \node[below] at (1,-.11) {$1$};
 \node[below] at (2,-.11) {$2$};
\end{tikzpicture}
\caption{Zbiór $[0,1)$ jest otwarty w $A=[0,2]$, a nie w $\mathbb R$.}
\label{fig:topog-podprzestrzen-przedzial}
\end{figure}

\begin{stwierdzenie}[Domkniętość i domknięcie względne]
\label{prop:topog-podprzestrzen-domkniecie}
Niech $A\subseteq X$, gdzie $X$ jest przestrzenią topologiczną.
Dla każdego $B\subseteq A$ zachodzą równoważność i równość:
\begin{enumerate}[label=(\roman*)]
 \item $B$ jest domknięty w $A$ wtedy i tylko wtedy, gdy
       $B=A\cap F$ dla pewnego zbioru domkniętego $F\subseteq X$;
 \item $\overline B^{\,A}=A\cap\overline B^{\,X}$.
\end{enumerate}
\end{stwierdzenie}
\begin{proof}
Zbiór $B$ jest domknięty w $A$ dokładnie wtedy, gdy
$A\setminus B=A\cap U$ dla pewnego otwartego $U\subseteq X$.
Biorąc dopełnienia w $A$, otrzymujemy
$B=A\cap(X\setminus U)$, co dowodzi (i).

Niech $C=A\cap\overline B^{\,X}$. Jest on domknięty
w $A$ i zawiera $B$, więc $\overline B^{\,A}\subseteq C$.
Jeśli $D\subseteq A$ jest dowolnym domkniętym podzbiorem $A$
zawierającym $B$, to na mocy (i) istnieje domknięty $F\subseteq X$
taki, że $D=A\cap F$. Z $B\subseteq F$ wynika
$\overline B^{\,X}\subseteq F$, a więc $C\subseteq D$.
Stąd $C$ jest najmniejszym domkniętym podzbiorem $A$
zawierającym $B$, co dowodzi (ii).
\end{proof}

\begin{stwierdzenie}[Wnętrze i brzeg względny]
\label{prop:topog-podprzestrzen-wnetrze}
Dla przestrzeni topologicznej $X$ oraz zbiorów $B\subseteq A\subseteq X$
mamy
\begin{align*}
 \operatorname{int}_A B
 &=A\setminus\overline{A\setminus B}^{\,X},\\
 \partial_A B
 &=A\cap\overline B^{\,X}\cap\overline{A\setminus B}^{\,X}.
\end{align*}
\end{stwierdzenie}
\begin{proof}
W dowolnej przestrzeni wnętrze zbioru jest dopełnieniem
domknięcia jego dopełnienia. Stosując tę własność w $A$ oraz
wzór na domknięcie względne, dostajemy
\[
\operatorname{int}_A B
 =A\setminus\overline{A\setminus B}^{\,A}
 =A\setminus(A\cap\overline{A\setminus B}^{\,X})
 =A\setminus\overline{A\setminus B}^{\,X}.
\]
Brzeg jest przecięciem domknięć zbioru i jego dopełnienia
w rozważanej przestrzeni; zastosowanie poprzedniego
stwierdzenia dwukrotnie daje drugą równość.
\end{proof}

\begin{uwaga}[Wnętrze względne nie jest przecięciem wnętrz]
W ogólności $\operatorname{int}_A B$ nie musi równać się
$A\cap\operatorname{int}_X B$. Dla $X=\mathbb R$,
$A=B=[0,1]$ mamy $\operatorname{int}_A B=A$, podczas gdy
$A\cap\operatorname{int}_{\mathbb R}B=(0,1)$.
\end{uwaga}

\begin{stwierdzenie}[Dziedziczenie topologii i włożenie]
\label{prop:topog-podprzestrzen-przechodniosc}
Niech $B\subseteq A\subseteq X$, gdzie $X$ jest przestrzenią
topologiczną. Topologia na $B$ odziedziczona najpierw z $X$
przez $A$, a następnie z $A$, jest równa topologii
odziedziczonej bezpośrednio z $X$. Ponadto kanoniczne
włożenie $\iota\colon A\to X$, $\iota(a)=a$, jest ciągłe.
\end{stwierdzenie}
\begin{proof}
Zbiór otwarty w $B$ otrzymany dwuetapowo ma postać
$B\cap(A\cap U)=B\cap U$, gdzie $U$ jest otwarty w $X$;
są to dokładnie zbiory topologii odziedziczonej bezpośrednio.
Dla otwartego $U\subseteq X$ przeciwobraz
$\iota^{-1}(U)=A\cap U$ jest otwarty w $A$, zatem
$\iota$ jest ciągłe.
\end{proof}

\begin{twierdzenie}[Ciągłość odwzorowania do podprzestrzeni]
\label{thm:topog-do-podprzestrzeni}
Niech $X,Y$ będą przestrzeniami topologicznymi,
$A\subseteq Y$, a $f\colon X\to A$ dowolnym odwzorowaniem.
Jest ono ciągłe do $A$ z topologią podprzestrzeni wtedy
i tylko wtedy, gdy złożenie $\iota\circ f\colon X\to Y$
z kanonicznym włożeniem $\iota\colon A\to Y$ jest ciągłe.
\end{twierdzenie}
\begin{proof}
Jeśli $f$ jest ciągłe, złożenie z ciągłym włożeniem
jest ciągłe. Odwrotnie, każdy otwarty $W\subseteq A$
ma postać $A\cap V$ dla otwartego $V\subseteq Y$.
Ponieważ wartości $f$ należą do $A$, zachodzi
$f^{-1}(W)=(\iota\circ f)^{-1}(V)$, a zatem
przeciwobraz $W$ jest otwarty.
\end{proof}

\section{Topologia iloczynu}
\label{sec:topog-iloczyn}

\begin{definicja}[Iloczyn przestrzeni topologicznych]
\label{def:topog-iloczyn}
Niech $I$ będzie zbiorem indeksów, a $(X_i,\tau_i)_{i\in I}$
rodziną przestrzeni topologicznych. Na zbiorze iloczynowym
$P=\prod_{i\in I}X_i$ rozważamy odwzorowania
$\pi_i\colon P\to X_i$, $\pi_i((x_j)_{j\in I})=x_i$.
\emph{Topologią iloczynu} na $P$ nazywamy topologię
generowaną przez rodzinę
\[
\mathcal S=\{\pi_i^{-1}(U):i\in I,\ U\in\tau_i\}.
\]
Jest ona zatem najmniejszą topologią, w której wszystkie
rzutowania $\pi_i$ są ciągłe. Dla $I=\varnothing$
przyjmujemy $P=\{()\}$, a topologia iloczynu jest
jedyną topologią na zbiorze jednoelementowym.
\end{definicja}

\begin{twierdzenie}[Baza topologii iloczynu]
\label{thm:topog-iloczyn-baza}
Niech $(X_i)_{i\in I}$ będzie rodziną przestrzeni
topologicznych, a $P=\prod_{i\in I}X_i$.
Bazę topologii iloczynu stanowią wszystkie zbiory
\[
\prod_{i\in I}U_i,\qquad
U_i\subseteq X_i\text{ otwarty},\qquad
U_i=X_i\text{ dla wszystkich poza skończenie wieloma }i.
\]
Dla $I=\varnothing$ zbiór bazowy jest równy $P$.
\end{twierdzenie}
\begin{proof}
Z definicji topologii generowanej przez podbazę jej
bazą są skończone przecięcia elementów $\mathcal S$,
włącznie z przecięciem pustym, które wynosi $P$.
Dla skończonego zbioru indeksów $F\subseteq I$ oraz
otwartych $U_i\subseteq X_i$ mamy
\[
\bigcap_{i\in F}\pi_i^{-1}(U_i)=
\prod_{i\in I}V_i,\qquad
V_i=\begin{cases}U_i,&i\in F,\\X_i,&i\notin F.\end{cases}
\]
Gdy ten sam indeks pojawia się kilkakrotnie, należy
zastąpić odpowiadające mu zbiory otwarte ich skończonym
przecięciem. Odwrotnie, każdy iloczyn z tezy jest takim
skończonym przecięciem przeciwobrazów.
\end{proof}

\begin{przyklad}[Prostokąty w $\mathbb R^2$]
\label{ex:topog-iloczyn-prostokaty}
Dla $X_1=X_2=\mathbb R$ z topologią euklidesową bazą
topologii iloczynu na $\mathbb R^2$ są prostokąty otwarte
$(a,b)\times(c,d)$. Topologia ta pokrywa się z topologią
metryki euklidesowej. Każdy punkt prostokąta ma kulę
euklidesową w nim zawartą, a każda kula euklidesowa
wokół danego punktu zawiera prostokąt otwarty zawierający
ten punkt. W pierwszą stronę wybieramy promień mniejszy
od odległości obu współrzędnych punktu do brzegów
odpowiednich przedziałów; w drugą wystarcza prostokąt
o półbokach $\varepsilon/3$ dla kuli o promieniu
$\varepsilon$, ponieważ odległość punktów w takim
prostokącie nie przekracza $\sqrt 2\,\varepsilon/3$.
\end{przyklad}

\begin{figure}[htbp]
\centering
\begin{tikzpicture}[>=Stealth,scale=.85]
 \draw[->] (-.3,0)--(5.3,0) node[right] {$X_1=\mathbb R$};
 \draw[->] (0,-.3)--(0,3.65) node[above] {$X_2=\mathbb R$};
 \fill[blue!13] (1.1,.7) rectangle (4.0,2.7);
 \draw[blue!70!black,thick,dashed] (1.1,.7) rectangle (4.0,2.7);
 \fill (2.6,1.75) circle (2.2pt);
 \node[above right] at (2.6,1.75) {$(x_1,x_2)$};
 \node[below] at (2.5,.7) {$U_1\times U_2$};
 \draw[blue!70!black,thick] (1.1,-.14)--(4,-.14);
 \draw[blue!70!black,thick] (-.14,.7)--(-.14,2.7);
 \node[below] at (1.1,-.14) {$a$};
 \node[below] at (4,-.14) {$b$};
 \node[left] at (-.14,.7) {$c$};
 \node[left] at (-.14,2.7) {$d$};
\end{tikzpicture}
\caption{Bazowe otoczenie punktu w topologii iloczynu dwóch prostych.}
\label{fig:topog-iloczyn-prostokat}
\end{figure}

\begin{twierdzenie}[Własność uniwersalna iloczynu]
\label{thm:topog-iloczyn-uniwersalnosc}
Niech $(X_i)_{i\in I}$ będzie rodziną przestrzeni
topologicznych, $P=\prod_{i\in I}X_i$ z topologią
iloczynu, a $Z$ dowolną przestrzenią topologiczną.
Odwzorowanie $f\colon Z\to P$ jest ciągłe wtedy i tylko
wtedy, gdy dla każdego $i\in I$ odwzorowanie
$\pi_i\circ f\colon Z\to X_i$ jest ciągłe.
\end{twierdzenie}
\begin{proof}
Jeśli $f$ jest ciągłe, to każde złożenie z ciągłym
rzutowaniem jest ciągłe. Odwrotnie, niech wszystkie
$\pi_i\circ f$ będą ciągłe. Dla dowolnego elementu
podbazy $\pi_i^{-1}(U)$ mamy
\[
f^{-1}(\pi_i^{-1}(U))=(\pi_i\circ f)^{-1}(U),
\]
a prawa strona jest otwarta w $Z$. Z kryterium ciągłości
przez podbazę wynika ciągłość $f$. Dotyczy to również
pustej rodziny indeksów, gdy $P$ jest jednoelementowe.
\end{proof}

\begin{stwierdzenie}[Rzutowania są odwzorowaniami otwartymi]
\label{prop:topog-rzutowania-otwarte}
Niech $P=\prod_{i\in I}X_i$ będzie iloczynem przestrzeni
topologicznych. Dla każdego $j\in I$ rzutowanie
$\pi_j\colon P\to X_j$ jest odwzorowaniem otwartym.
\end{stwierdzenie}
\begin{proof}
Niech $W=\prod_{i\in I}U_i$ będzie bazowym zbiorem
otwartym, z $U_i=X_i$ poza skończonym zbiorem indeksów.
Jeśli $W=\varnothing$, to $\pi_j(W)=\varnothing$ jest
otwarty. W przeciwnym przypadku istnieje punkt
$(x_i)_{i\in I}\in W$. Dla dowolnego $y\in U_j$
rodzina, która ma współrzędną $y$ w miejscu $j$, a
pozostałe współrzędne $x_i$, należy do $W$. Zatem
$\pi_j(W)=U_j$, zbiór otwarty. Każdy zbiór otwarty
w $P$ jest sumą takich zbiorów bazowych, a obraz
zachowuje dowolne sumy. Stąd rzutowanie jest otwarte.
\end{proof}

\begin{stwierdzenie}[Domknięcie prostokąta]
\label{prop:topog-iloczyn-domkniecie}
Niech $X,Y$ będą przestrzeniami topologicznymi oraz
$A\subseteq X$, $B\subseteq Y$. W topologii iloczynu
na $X\times Y$ zachodzi
\[
\overline{A\times B}^{\,X\times Y}
 =\overline A^{\,X}\times\overline B^{\,Y}.
\]
\end{stwierdzenie}
\begin{proof}
Jeśli $A=\varnothing$ lub $B=\varnothing$, obie strony
są puste. W przeciwnym przypadku punkt $(x,y)$ należy
do domknięcia $A\times B$ wtedy i tylko wtedy, gdy każdy
bazowy prostokąt $U\times V$ zawierający $(x,y)$ przecina
$A\times B$. To jest równoważne temu, że każde otwarte
otoczenie $U$ punktu $x$ przecina $A$ i każde otwarte
otoczenie $V$ punktu $y$ przecina $B$.
Z definicji domknięcia dostajemy
$x\in\overline A^{\,X}$ i $y\in\overline B^{\,Y}$.
\end{proof}

\begin{uwaga}[Dlaczego skończenie wiele współrzędnych?]
\label{rem:topog-iloczyn-skrzynkowa}
Przy nieskończonym zbiorze indeksów żądanie, aby
\emph{wszystkie} iloczyny $\prod_i U_i$ z otwartymi
$U_i\subseteq X_i$ były otwarte, może dać inną,
\emph{topologię skrzynkową}. Dla $X_n=\mathbb R$ i
$I=\mathbb N$ zbiór $\prod_{n\ge1}(-1/n,1/n)$ jest
otwarty w topologii skrzynkowej, lecz nie w topologii
iloczynu: każde bazowe otoczenie punktu zerowego
pozostawia wszystkie poza skończenie wieloma współrzędne
bez ograniczeń, więc nie może zawierać się w podanym
zbiorze. Dla skończonej liczby czynników obie topologie
są równe.
\end{uwaga}

\section{Suma rozłączna przestrzeni}
\label{sec:topog-suma-rozlaczna}

\begin{definicja}[Suma rozłączna]
\label{def:topog-suma-rozlaczna}
Niech $(X_i,\tau_i)_{i\in I}$ będzie rodziną przestrzeni
topologicznych. \emph{Sumą rozłączną} jest zbiór
$S=\coprod_{i\in I}X_i=\{(i,x):i\in I,\ x\in X_i\}$
z topologią określoną następująco: $W\subseteq S$ jest
otwarty wtedy i tylko wtedy, gdy dla każdego $i\in I$
zbiór $W_i=\{x\in X_i:(i,x)\in W\}$ jest otwarty
w $X_i$. Naturalne włożenia oznaczamy przez
$\jmath_i\colon X_i\to S$, $\jmath_i(x)=(i,x)$.
Dla $I=\varnothing$ otrzymujemy pustą przestrzeń.
\end{definicja}

\begin{stwierdzenie}[Topologia i składniki sumy rozłącznej]
\label{prop:topog-suma-rozlaczna-topologia}
Definicja \ref{def:topog-suma-rozlaczna} wyznacza topologię
na $S$. Dla każdego $i\in I$ odwzorowanie $\jmath_i$
jest zanurzeniem topologicznym, a jego obraz
$\jmath_i(X_i)$ jest zarówno otwarty, jak i domknięty w $S$.
\end{stwierdzenie}
\begin{proof}
Przekroje zbiorów pustego i całego $S$ są odpowiednio
puste i równe $X_i$. Ponadto przekrój sumy dowolnej
rodziny zbiorów jest sumą ich przekrojów, a przekrój
przecięcia dwóch zbiorów jest przecięciem ich przekrojów.
Aksjomaty topologii wynikają z aksjomatów w każdym $X_i$.

Zbiór $\jmath_i(X_i)$ jest otwarty, ponieważ jego przekrój
w składniku $i$ jest całym $X_i$, a w każdym innym jest
pusty. Jego dopełnienie także jest otwarte, zatem obraz
jest domknięty. Jeśli $U\subseteq X_i$ jest otwarty,
to $\jmath_i(U)$ jest otwarty w $S$, gdyż tylko przekrój
w składniku $i$ jest równy $U$, a pozostałe są puste.
Dla otwartego $W\subseteq S$ przeciwobraz
$\jmath_i^{-1}(W)=W_i$ jest otwarty w $X_i$.
Odwzorowanie $\jmath_i\colon X_i\to\jmath_i(X_i)$
jest więc ciągłą i otwartą bijekcją, czyli homeomorfizmem.
\end{proof}

\begin{twierdzenie}[Własność uniwersalna sumy rozłącznej]
\label{thm:topog-suma-uniwersalnosc}
Niech $S=\coprod_{i\in I}X_i$ będzie sumą rozłączną
przestrzeni topologicznych, $Y$ przestrzenią topologiczną,
a $f\colon S\to Y$ dowolnym odwzorowaniem. Wtedy $f$
jest ciągłe wtedy i tylko wtedy, gdy dla każdego $i\in I$
odwzorowanie $f\circ\jmath_i\colon X_i\to Y$ jest ciągłe.
\end{twierdzenie}
\begin{proof}
Jeśli $f$ jest ciągłe, to wszystkie złożenia z ciągłymi
$\jmath_i$ są ciągłe. Odwrotnie, niech wszystkie złożenia
będą ciągłe. Dla otwartego $V\subseteq Y$ oraz każdego
$i\in I$ mamy
$\jmath_i^{-1}(f^{-1}(V))=(f\circ\jmath_i)^{-1}(V)$,
zatem wszystkie przekroje $f^{-1}(V)$ są otwarte.
Z definicji topologii na $S$ zbiór $f^{-1}(V)$ jest otwarty.
\end{proof}

\begin{przyklad}[Dwie kopie tej samej przestrzeni]
\label{ex:topog-dwie-kopie}
Przestrzeń $\mathbb R\coprod\mathbb R$ składa się z dwóch
rozłącznych kopii prostej. Każda kopia jest otwarta i
domknięta. Nie jest to zwykła suma zbiorów
$\mathbb R\cup\mathbb R=\mathbb R$: znaczniki składników
odróżniają punkty o tej samej współrzędnej.
Nie oznacza to, że suma rozłączna nie może być homeomorficzna
z pewną podprzestrzenią prostej; chodzi o różnicę
między samymi konstrukcjami zbiorów.
Rozłączność początkowych zbiorów $X_i$ nie jest wymagana.
\end{przyklad}

\section{Topologie początkowe i końcowe}
\label{sec:topog-topologie-poczatkowe-koncowe}

\begin{definicja}[Topologia początkowa]
\label{def:topog-topologia-poczatkowa}
Niech $X$ będzie zbiorem, $(Y_i,\tau_i)_{i\in I}$ rodziną
przestrzeni topologicznych, a $f_i\colon X\to Y_i$
odwzorowaniami. \emph{Topologia początkowa} wyznaczona
przez rodzinę $(f_i)_{i\in I}$ jest topologią na $X$
generowaną przez podbazę
\[
\mathcal S=\{f_i^{-1}(U):i\in I,\ U\in\tau_i\}.
\]
Jest to najmniejsza topologia na $X$, dla której wszystkie
$f_i$ są ciągłe. Jeśli $I=\varnothing$, jest to
topologia antydyskretna $\{\varnothing,X\}$,
rozumiana jako jednoelementowa, gdy $X=\varnothing$.
\end{definicja}

\begin{twierdzenie}[Własność uniwersalna topologii początkowej]
\label{thm:topog-poczatkowa-uniwersalnosc}
Niech $X$ ma topologię początkową wyznaczoną przez
$f_i\colon X\to Y_i$, $i\in I$, a $Z$ będzie dowolną
przestrzenią topologiczną. Odwzorowanie $g\colon Z\to X$
jest ciągłe wtedy i tylko wtedy, gdy wszystkie złożenia
$f_i\circ g\colon Z\to Y_i$ są ciągłe.
\end{twierdzenie}
\begin{proof}
Jeden kierunek wynika ze składania odwzorowań ciągłych.
Jeśli wszystkie złożenia są ciągłe, to dla $i\in I$
i otwartego $U\subseteq Y_i$ zachodzi
$g^{-1}(f_i^{-1}(U))=(f_i\circ g)^{-1}(U)$,
a prawa strona jest otwarta w $Z$.
Ciągłość $g$ wynika z kryterium przez podbazę.
\end{proof}

\begin{uwaga}[Podprzestrzeń i iloczyn jako konstrukcje początkowe]
Topologia podprzestrzeni na $A\subseteq X$ jest początkowa
względem włożenia $\iota\colon A\to X$. Topologia iloczynu
na $\prod_i X_i$ jest początkowa względem rzutowań $\pi_i$.
Oba przypadki wynikają bezpośrednio z definicji rodzin
podbazowych. Ich własności uniwersalne są więc
szczególnymi przypadkami twierdzenia
\ref{thm:topog-poczatkowa-uniwersalnosc}.
\end{uwaga}

\begin{definicja}[Topologia końcowa]
\label{def:topog-topologia-koncowa}
Niech $(X_i,\tau_i)_{i\in I}$ będzie rodziną przestrzeni
topologicznych, $Y$ zbiorem oraz
$f_i\colon X_i\to Y$ odwzorowaniami. \emph{Topologią końcową}
wyznaczoną przez rodzinę $(f_i)_{i\in I}$ nazywamy
\[
\tau_{\mathrm k}=\{V\subseteq Y:
 f_i^{-1}(V)\in\tau_i\text{ dla każdego }i\in I\}.
\]
Jest to największa topologia na $Y$, dla której wszystkie
$f_i$ są ciągłe. Jeśli $I=\varnothing$, otrzymujemy
topologię dyskretną $\mathcal P(Y)$.
\end{definicja}

\begin{stwierdzenie}[Poprawność definicji topologii końcowej]
\label{prop:topog-koncowa-topologia}
Rodzina $\tau_{\mathrm k}$ z definicji
\ref{def:topog-topologia-koncowa} jest topologią.
Każda topologia na $Y$, dla której wszystkie $f_i$
są ciągłe, jest zawarta w $\tau_{\mathrm k}$.
\end{stwierdzenie}
\begin{proof}
Przeciwobrazy $\varnothing$ i $Y$ są odpowiednio
$\varnothing$ i $X_i$, a więc są otwarte.
Przeciwobrazy zachowują dowolne sumy i skończone
przecięcia zbiorów, co dowodzi aksjomatów topologii.
Jeśli w pewnej topologii na $Y$ wszystkie $f_i$
są ciągłe, każdy jej zbiór otwarty ma otwarte
przeciwobrazy przy wszystkich $f_i$, należy więc
do $\tau_{\mathrm k}$.
\end{proof}

\begin{twierdzenie}[Własność uniwersalna topologii końcowej]
\label{thm:topog-koncowa-uniwersalnosc}
Niech na zbiorze $Y$ zadana będzie topologia końcowa
względem $f_i\colon X_i\to Y$, $i\in I$, i niech
$Z$ będzie przestrzenią topologiczną. Odwzorowanie
$h\colon Y\to Z$ jest ciągłe wtedy i tylko wtedy, gdy
wszystkie złożenia $h\circ f_i\colon X_i\to Z$
są ciągłe.
\end{twierdzenie}
\begin{proof}
Jeśli $h$ jest ciągłe, ciągłość złożeń wynika z ciągłości
$f_i$. Jeśli wszystkie złożenia są ciągłe, to dla dowolnego
otwartego $W\subseteq Z$ wszystkie zbiory
\[
f_i^{-1}(h^{-1}(W))=(h\circ f_i)^{-1}(W)
\]
są otwarte. Na mocy definicji topologii końcowej
$h^{-1}(W)$ jest otwarty w $Y$.
\end{proof}

\begin{uwaga}[Suma rozłączna jako konstrukcja końcowa]
Topologia na $\coprod_i X_i$ jest końcowa względem
włożeń $\jmath_i$. W istocie warunek otwartości
$\jmath_i^{-1}(W)$ jest dokładnie warunkiem otwartości
odpowiedniego przekroju $W_i$ z definicji sumy rozłącznej.
\end{uwaga}

\section{Topologia ilorazowa}
\label{sec:topog-iloraz}

\begin{definicja}[Przestrzeń ilorazowa]
\label{def:topog-przestrzen-ilorazowa}
Niech $(X,\tau)$ będzie przestrzenią topologiczną,
a $\sim$ relacją równoważności na $X$. Niech $X/{\sim}$
oznacza zbiór klas równoważności, a
$q\colon X\to X/{\sim}$, $q(x)=[x]$, kanoniczną surjekcję.
\emph{Topologią ilorazową} jest rodzina
\[
\tau_{/{\sim}}=\{V\subseteq X/{\sim}:q^{-1}(V)\in\tau\}.
\]
Przestrzeń $(X/{\sim},\tau_{/{\sim}})$ nazywamy
\emph{przestrzenią ilorazową}. Topologia ilorazowa jest
topologią końcową względem $q$; w szczególności $q$ jest ciągłe.
\end{definicja}

\begin{definicja}[Odwzorowanie ilorazowe i zbiory nasycone]
\label{def:topog-odwzorowanie-ilorazowe}
Niech $q\colon X\to Y$ będzie ciągłą surjekcją
przestrzeni topologicznych. Nazywamy ją
\emph{odwzorowaniem ilorazowym}, jeśli
$V\subseteq Y$ jest otwarty dokładnie wtedy, gdy
$q^{-1}(V)$ jest otwarty w $X$.
Zbiór $A\subseteq X$ nazywamy \emph{nasyconym względem $q$},
jeśli $q^{-1}(q(A))=A$. Jest to równoważne temu,
że $A$ jest sumą włókien $q^{-1}(y)$.
\end{definicja}

\begin{stwierdzenie}[Zbiory nasycone i otwartość w ilorazie]
\label{prop:topog-iloraz-nasycone}
Niech $q\colon X\to Y$ będzie odwzorowaniem ilorazowym.
Dla dowolnego zbioru $A\subseteq X$ zachodzi:
\begin{enumerate}[label=(\roman*)]
\item $q(A)$ jest otwarty w $Y$ wtedy i tylko wtedy, gdy
      $q^{-1}(q(A))$ jest otwarty w $X$;
\item $q(A)$ jest domknięty w $Y$ wtedy i tylko wtedy, gdy
      $q^{-1}(q(A))$ jest domknięty w $X$.
\end{enumerate}
Jeśli $A$ jest nasycony, warunki po prawej stronach
sprowadzają się do otwartości lub domkniętości $A$.
\end{stwierdzenie}
\begin{proof}
Pierwsza równoważność jest definicją odwzorowania
ilorazowego zastosowaną do zbioru $q(A)$.
Dla drugiej stosujemy tę samą definicję do dopełnienia
$Y\setminus q(A)$ i używamy równości
$q^{-1}(Y\setminus q(A))=X\setminus q^{-1}(q(A))$.
Ostatnie zdanie wynika z definicji nasycenia.
\end{proof}

\begin{twierdzenie}[Własność uniwersalna ilorazu]
\label{thm:topog-iloraz-uniwersalnosc}
Niech $q\colon X\to Y$ będzie odwzorowaniem ilorazowym,
$Z$ dowolną przestrzenią topologiczną, a $h\colon Y\to Z$
odwzorowaniem. Wtedy $h$ jest ciągłe wtedy i tylko wtedy,
gdy $h\circ q\colon X\to Z$ jest ciągłe.
Ponadto jeśli $g\colon X\to Z$ jest stałe na każdym
włóknie $q$, istnieje dokładnie jedno odwzorowanie
$h\colon Y\to Z$ z $g=h\circ q$, a $h$ jest ciągłe
dokładnie wtedy, gdy $g$ jest ciągłe.
\end{twierdzenie}
\begin{proof}
Pierwsza równoważność jest własnością uniwersalną
topologii końcowej względem jednego odwzorowania.
Jeśli $g$ ma stałą wartość na każdym włóknie, określamy
$h(y)=g(x)$, gdzie $x$ jest dowolnym punktem z
$q^{-1}(y)$. Włókno jest niepuste z surjektywności,
a stałość $g$ zapewnia niezależność od wyboru $x$.
Wówczas $h(q(x))=g(x)$ dla wszystkich $x\in X$.
Surjektywność $q$ gwarantuje jednoznaczność $h$.
Równoważność ciągłości wynika z pierwszej części.
\end{proof}

\begin{stwierdzenie}[Wystarczające kryterium ilorazowości]
\label{prop:topog-otwarta-surjekcja-iloraz}
Każda ciągła surjekcja $q\colon X\to Y$, która jest
odwzorowaniem otwartym lub odwzorowaniem domkniętym,
jest odwzorowaniem ilorazowym.
\end{stwierdzenie}
\begin{proof}
Ciągłość daje: dla każdego otwartego $V\subseteq Y$
przeciwobraz $q^{-1}(V)$ jest otwarty w $X$.
Jeśli $q$ jest otwarte i $q^{-1}(V)$ jest otwarty,
to $V=q(q^{-1}(V))$ jest otwarty dzięki surjektywności.
Jeśli $q$ jest domknięte i $q^{-1}(V)$ jest otwarty,
to $q^{-1}(Y\setminus V)=X\setminus q^{-1}(V)$
jest domknięty; zatem $Y\setminus V$ jest domknięty
jako obraz zbioru domkniętego. W obu przypadkach
otrzymujemy równoważność definiującą ilorazowość.
\end{proof}

\begin{przyklad}[Sklejenie dwóch punktów prostej]
\label{ex:topog-iloraz-dwa-punkty}
W $\mathbb R$ z topologią euklidesową przyjmijmy
$x\sim y$ wtedy i tylko wtedy, gdy $x=y$ albo
$\{x,y\}=\{0,1\}$. Iloraz otrzymujemy, utożsamiając
$0$ i $1$, a pozostałe punkty pozostają rozróżnione.
Dla $U=(-1/2,1/2)$ zbiór nasycony jest równy
$q^{-1}(q(U))=U\cup\{1\}$, który nie jest otwarty
w $\mathbb R$. Wobec tego $q(U)$ nie jest otwarty
w ilorazie, mimo że $U$ jest otwarty w $\mathbb R$.
Zatem odwzorowanie ilorazowe nie musi być otwarte.
\end{przyklad}

\begin{przyklad}[Iloraz nie musi być odwzorowaniem domkniętym]
\label{ex:topog-iloraz-niedomkniety}
Na $\mathbb R$ utożsammy ze sobą wszystkie punkty
przedziału $(0,1)$, nie utożsamiając żadnych innych.
Niech $q$ będzie odwzorowaniem do ilorazu.
Zbiór $F=\{1/2\}$ jest domknięty w $\mathbb R$,
lecz $q^{-1}(q(F))=(0,1)$ nie jest domknięty.
Z kryterium zbiorów nasyconych wynika, że $q(F)$
nie jest domknięty. Odwzorowanie ilorazowe nie musi
więc być domknięte.
\end{przyklad}

\begin{przyklad}[Ciągła surjekcja, która nie jest ilorazowa]
\label{ex:topog-surjekcja-nieilorazowa}
Niech $S=\{0,1\}$ i rozważmy odwzorowanie
identycznościowe
\[
\operatorname{id}\colon(S,\mathcal P(S))\longrightarrow
       (S,\{\varnothing,S\}).
\]
Jest ono ciągłą surjekcją, ponieważ przeciwdziedzina
ma topologię antydyskretną. Nie jest ilorazowe:
przeciwobraz $\{0\}$ jest otwarty w dyskretnej dziedzinie,
ale $\{0\}$ nie jest otwarty w przeciwdziedzinie.
\end{przyklad}

\section{Zestawienie własności uniwersalnych}
\label{sec:topog-konstrukcje-podsumowanie}

Wszystkie konstrukcje rozdziału powstają z tego samego
pytania: jakie zbiory uznać za otwarte, aby wskazane
odwzorowania były ciągłe?
Dla topologii początkowej wymagamy ciągłości odwzorowań
\emph{do} danych przestrzeni i wybieramy najmniejszą
taką topologię. Dla topologii końcowej wymagamy ciągłości
odwzorowań \emph{z} danych przestrzeni i wybieramy
największą taką topologię. Iloczyny i podprzestrzenie
są szczególnymi przypadkami pierwszej konstrukcji;
sumy rozłączne i ilorazy --- drugiej.

\begin{uwaga}[Kierunek sprawdzania ciągłości]
W topologii podprzestrzeni i iloczynu sprawdzamy ciągłość
odwzorowań $g\colon Z\to X$ na podstawie ich złożeń
z włożeniami lub rzutowaniami \emph{wychodzącymi z $X$}.
W sumie rozłącznej i topologii ilorazowej sprawdzamy
ciągłość odwzorowań $h\colon Y\to Z$ na podstawie
złożeń z naturalnymi odwzorowaniami \emph{wchodzącymi do $Y$}.
Rozróżnienie to będzie wykorzystywane przy konstrukcjach
bardziej złożonych przestrzeni i badaniu ich niezmienników.
\end{uwaga}